% =====================================================================
% ejercicio03.tex
% Responsabilidad única: relacionar la sigmoide con los log-odds,
% clasificar tres máquinas bajo distintos umbrales, deducir la frontera
% de decisión y el umbral que minimiza el costo esperado. Cifras
% verificadas con verificacion/teoria.py.
% =====================================================================
\subsection{Ejercicio 3 --- Regresión logística: probabilidad, odds y umbral}

El modelo estima la probabilidad de falla en la próxima semana ($Y=1$)
a partir de los años de uso $x_1$ y las alertas del último mes $x_2$:
\[
    z = -3 + 0.6x_1 + 1.2x_2, \qquad p = \sigma(z) = \frac{1}{1+e^{-z}} .
\]

\begin{enumerate}

\item \textbf{Log-odds y factor de odds.}
A partir de la definición de $p$,
\[
    1-p = 1 - \frac{1}{1+e^{-z}} = \frac{e^{-z}}{1+e^{-z}}
    \;\Longrightarrow\;
    \frac{p}{1-p} = \frac{1}{e^{-z}} = e^{z}
    \;\Longrightarrow\;
    \log\frac{p}{1-p} = z ,
\]
es decir, el modelo es \emph{lineal en los log-odds}, no en la
probabilidad. Si $x_1$ aumenta en una unidad con $x_2$ fijo, $z$ aumenta
exactamente $0.6$ y los odds pasan de $e^{z}$ a $e^{z+0.6}$, de forma que
se multiplican por
\[
    e^{0.6} \approx 1.8221 ,
\]
lo cual equivale a un aumento del 82.21\,\% \emph{en los odds},
constante sin importar el punto de partida. Esto no debe confundirse con
un cambio en puntos porcentuales de la probabilidad, que sí depende del
punto de partida porque la sigmoide no es lineal: pasar de $(2,1)$ a
$(3,1)$ lleva $p$ de $0.3543$ a $0.5000$ ($+14.57$ puntos porcentuales),
mientras que pasar de $(4,2)$ a $(5,2)$ lleva $p$ de $0.8581$ a $0.9168$
(apenas $+5.87$ puntos), aun cuando en ambos casos los odds se
multiplican por el mismo $1.8221$.

\item \textbf{Probabilidades y clases para $A$, $B$, $C$.}
Usando $\hat y=1$ cuando $p\ge t$:

\begin{table}[H]
\centering
\begin{tabular}{crrrrcc}
\toprule
Máquina & $(x_1,x_2)$ & $z$ & $p=\sigma(z)$ & odds & $\hat y$ ($t=0.5$) & $\hat y$ ($t=0.8$) \\
\midrule
$A$ & $(2,1)$ & $-3+1.2+1.2=-0.6$ & 0.3543 & 0.5488 & 0 & 0 \\
$B$ & $(5,0)$ & $-3+3.0+0\phantom{.0}=\phantom{-}0.0$ & 0.5000 & 1.0000 & 1 & 0 \\
$C$ & $(4,2)$ & $-3+2.4+2.4=\phantom{-}1.8$ & 0.8581 & 6.0496 & 1 & 1 \\
\bottomrule
\end{tabular}
\caption{Puntaje lineal, probabilidad de falla y clase predicha por umbral.}
\label{tab:ej3-maquinas}
\end{table}

Cabe mencionar el caso de $B$, que cae exactamente sobre la frontera de
$t=0.5$ ($p=0.5$): con la regla $p\ge t$ se clasifica como falla, pero
con una regla estricta $p>t$ se clasificaría como no falla, por lo cual
la convención de desempate debe declararse explícitamente. Al subir el
umbral a $0.8$ solo la máquina $C$ conserva la predicción de falla.

\item \textbf{Frontera de decisión para un umbral general.}
Como $\sigma$ es estrictamente creciente, $p\ge t$ equivale a
$z\ge\sigma^{-1}(t)=\log\frac{t}{1-t}$, por lo que la frontera es la
recta
\[
    0.6x_1 + 1.2x_2 = 3 + \log\frac{t}{1-t}
    \quad\Longleftrightarrow\quad
    x_2 = \frac{3+\log\frac{t}{1-t}}{1.2} - 0.5\,x_1 ,
\]
y se predice falla en el semiplano por encima de ella. Particularizando,
\[
    \begin{aligned}
    t=0.5:&\quad 0.6x_1+1.2x_2 = 3 &&\Longleftrightarrow\quad x_2 = 2.5 - 0.5x_1,\\
    t=0.8:&\quad 0.6x_1+1.2x_2 = 3+\log 4 \approx 4.3863 &&\Longleftrightarrow\quad x_2 \approx 3.6552 - 0.5x_1.
    \end{aligned}
\]
Al modificar solamente el umbral \textbf{la frontera se desplaza
paralelamente} (cambia el término independiente, de forma que con
$t=0.8$ se exige más evidencia y la región de falla se reduce), pero su
pendiente $-0.5=-\beta_1/\beta_2$ y su orientación quedan iguales, así
como \textbf{las probabilidades no cambian en absoluto}: $p(x)$ depende
solo de $\beta$, de modo que el umbral es una decisión posterior sobre
cómo convertir probabilidades en clases, no una propiedad del modelo
ajustado.

\begin{figure}[H]
\centering
\begin{tikzpicture}
\begin{axis}[
    width=0.66\textwidth, height=6cm,
    xmin=0, xmax=7.5, ymin=0, ymax=4.2,
    xlabel={$x_1$ (años de uso)}, ylabel={$x_2$ (alertas)},
    legend pos=north east, legend cell align=left,
    legend style={font=\small}, tick label style={font=\small},
]
\addplot[domain=0:7.5, thick, color=classblue] {2.5 - 0.5*x};
\addlegendentry{$t=0.5$}
\addplot[domain=0:7.5, thick, dashed, color=classred] {3.6552 - 0.5*x};
\addlegendentry{$t=0.8$}
\addplot[domain=0:7.5, thick, dotted, color=classgreen] {1.3448 - 0.5*x};
\addlegendentry{$t=0.2$ (inciso d)}
\addplot[only marks, mark=*, mark size=2.4pt, color=black,
         nodes near coords, point meta=explicit symbolic,
         every node near coord/.append style={font=\small, anchor=south west}]
    coordinates {(2,1)[$A$] (5,0)[$B$] (4,2)[$C$]};
\end{axis}
\end{tikzpicture}
\caption{Fronteras de decisión $p=t$ en el plano $(x_1,x_2)$: rectas
paralelas de pendiente $-0.5$; por encima de cada recta se predice falla.}
\end{figure}

\item \textbf{Umbral de costo mínimo.}
Con probabilidades calibradas, si se predice $0$ el único error posible
es un falso negativo (con probabilidad $p$ y costo $8$), y si se predice
$1$ el único error posible es un falso positivo (con probabilidad $1-p$
y costo $2$), de forma que
\[
    \mathbb{E}[\text{costo}\mid\hat y=0] = 8p,
    \qquad
    \mathbb{E}[\text{costo}\mid\hat y=1] = 2(1-p),
\]
y conviene predecir $1$ cuando $2(1-p)\le8p$, es decir, cuando
$p\ge 2/(2+8)=\mathbf{0.2}$ (en general,
$t^*=c_{FP}/(c_{FP}+c_{FN})$). Como un falso negativo cuesta
cuatro veces más que un falso positivo, el umbral óptimo baja de $0.5$
a $0.2$, y aplicándolo, $A$ ($p=0.3543$), $B$ ($p=0.5$) y $C$
($p=0.8581$) superan $0.2$, por lo cual \textbf{las tres máquinas se
clasifican como falla}, incluida $A$, que con $t=0.5$ no se marcaba
(su frontera aparece punteada en la figura anterior,
$x_2\approx1.3448-0.5x_1$).

\item \textbf{Mínimos cuadrados vs.\ sigmoide en $[0,1]$.}
Un modelo $\hat y=\beta_0+\beta^\top x$ es una función afín no acotada:
si algún $\beta_j\neq0$, al crecer o decrecer $x_j$ lo suficiente
$\hat y$ supera $1$ o cae por debajo de $0$, y mínimos cuadrados no
incluye ninguna restricción que lo evite, pues solo minimiza la
distancia cuadrática a etiquetas $0/1$ (por ejemplo, ajustando MCO a
$x=(0,1,2,3,10)$ con $y=(0,0,1,1,1)$ se obtiene
$\hat y\approx0.3248+0.0860x$, que predice $\hat y(15)\approx1.6146$, un
valor imposible como probabilidad). En cambio, la sigmoide cumple
$\sigma(z)=1/(1+e^{-z})\in(0,1)$ para todo $z\in\mathbb{R}$, porque
$e^{-z}>0$ hace que el denominador sea siempre mayor que $1$, de modo
que $\sigma$ transforma el puntaje lineal no acotado en una
probabilidad válida, tendiendo a $0$ y a $1$ solo asintóticamente.

\end{enumerate}
